\section*{Chapitre \ref{ch:b3:innate-immunity} --- Immunité innée et inflammation}

\begin{solution}{exo:b3:innate-immunity:1}
\omterm{def:b3:innate-immunity:innate}{Barrières}~: la kératine de la peau (mécanique), le mucus et les cils des
voies aériennes (piégeage et évacuation), l'acidité de l'estomac (tue
les microbes avalés), le lysozyme et les défensines des sécrétions
(lysent et perforent), le \omterm{def:b3:microbiomes:symbiosis}{microbiote} commensal (occupe les niches).
Cellules~: \omterm{def:b3:innate-immunity:innate}{neutrophiles} (tuent les bactéries par \omterm{def:b3:innate-immunity:killing}{phagocytose}),
\omterm{def:b3:innate-immunity:innate}{macrophages} (mangent, nettoient, donnent l'alerte), \omterm{def:b3:innate-immunity:innate}{cellules
dendritiques} (portent l'antigène aux ganglions et démarrent les réponses
adaptatives), mastocytes (libèrent l'histamine, enflamment), \omterm{def:b3:innate-immunity:innate}{cellules
tueuses naturelles} (tuent les cellules infectées et transformées).
\end{solution}

\begin{solution}{exo:b3:innate-immunity:2}
Partagé par toute une classe de microbes, essentiel à eux (il ne peut
donc être perdu pour s'échapper) et absent de l'hôte. \omterm{def:b3:bacteriology:envelope}{Lipopolysaccharide}
--- TLR4~; flagelline --- TLR5~; ARN double brin --- TLR3 (et RIG-I dans
le cytosol)~; ADN \omterm{def:b3:chromatin-epigenetics:methylation}{CpG} non méthylé --- TLR9~; fragments de \omterm{def:b3:bacteriology:envelope}{peptidoglycane}
--- protéines NOD~; $\beta$-glucane --- dectine-1.
\end{solution}

\begin{solution}{exo:b3:innate-immunity:3}
Rougeur~: dilatation des artérioles par l'histamine, le monoxyde d'azote
et les prostaglandines. Chaleur~: le même surcroît de débit sanguin.
Gonflement~: fuite plasmatique des veinules après contraction
endothéliale (histamine, C3a, C5a). Douleur~: la bradykinine et la
prostaglandine E$_{2}$ sensibilisent les nocicepteurs, plus la pression
du gonflement.
\end{solution}

\begin{solution}{exo:b3:innate-immunity:4}
Classique~: C1q liant un anticorps (ou la \omterm{prop:b3:innate-immunity:systemic}{protéine C réactive}) sur une
surface. Lectines~: la lectine liant le mannose sur les sucres
microbiens. Alterne~: bruit de fond spontané du C3 et dépôt de C3b sur
toute surface non protégée. Conséquences~: le C3b opsonise pour la
\omterm{def:b3:innate-immunity:killing}{phagocytose}~; C3a et C5a enflamment et recrutent~; C5b démarre le
\omterm{def:b3:innate-immunity:complement}{complexe d'attaque membranaire} qui lyse.
\end{solution}

\begin{solution}{exo:b3:innate-immunity:5}
$B = [k_{0}/(a-d)](e^{(a-d)t} - 1)$ avec $a - d = 0.1$~: à \qty{60}{s},
$20(e^{6} - 1) \approx 8000$~; à \qty{90}{s}, $20(e^{9} - 1) \approx
1.6\times 10^{5}$. Cellule de l'hôte~: $k_{0}/(d - a) = 2/0.1 = 20$.
\end{solution}

\begin{solution}{exo:b3:innate-immunity:6}
Roulement sur les \omterm{def:b3:innate-immunity:inflammation}{sélectines} (lectines liant un sucre)~; signalisation
par \omterm{def:b3:innate-immunity:inflammation}{chimiokine} qui active les \omterm{def:b3:innate-immunity:inflammation}{intégrines}~; adhérence ferme des
\omterm{def:b3:innate-immunity:inflammation}{intégrines} aux molécules d'adhérence endothéliales (ICAM)~;
transmigration entre les cellules endothéliales~; chimiotactisme le long
du gradient. Sans \omterm{def:b3:innate-immunity:inflammation}{intégrines}
$\beta_{2}$ les \omterm{def:b3:innate-immunity:innate}{neutrophiles} roulent mais ne s'arrêtent ni ne traversent
jamais~: ils s'accumulent dans le sang (comptes très élevés), les
infections récidivent sans former de pus, les plaies cicatrisent mal et
le cordon ombilical se détache tard.
\end{solution}

\begin{solution}{exo:b3:innate-immunity:7}
Granulomatose septique chronique~: la \omterm{def:b3:innate-immunity:killing}{NADPH oxydase} est défectueuse, si
bien qu'il ne se fait dans le phagosome ni superoxyde, ni peroxyde
d'hydrogène, ni hypochlorite. Infections par des organismes qui
détruisent leur propre peroxyde d'hydrogène (catalase positive)~:
\emph{Staphylococcus aureus}, \emph{Aspergillus},
\emph{Burkholderia}, \emph{Serratia}, \emph{Nocardia}. Des granulomes se
forment parce que les \omterm{def:b3:innate-immunity:innate}{macrophages} englobent des microbes qu'ils ne
peuvent tuer et les cloisonnent, avec un recrutement continuel, dans une
lésion chronique.
\end{solution}

\begin{solution}{exo:b3:innate-immunity:8}
Les récepteurs inhibiteurs comptent le CMH~I~; les récepteurs
activateurs comptent les ligands de stress~; la cellule meurt si
l'activation l'emporte sur l'inhibition. (a) Perdre le CMH~I ôte
l'inhibition~: la \omterm{def:b3:innate-immunity:innate}{cellule tueuse naturelle} la tue --- le prix à payer
pour échapper aux lymphocytes~T. (b) Avec le CMH~I conservé mais des
ligands de stress exposés, l'activation peut encore l'emporter et la
cellule est tuée~; les \omterm{def:b3:cancer-biology:hallmarks}{tumeurs} perdent donc souvent, ou répriment, les
ligands de stress.
\end{solution}

\begin{solution}{exo:b3:innate-immunity:9}
Une protéine pure ne porte aucun motif de pathogène~: les \omterm{def:b3:innate-immunity:innate}{cellules
dendritiques} qui la captent ne mûrissent donc pas, la présentent sans
costimulation, et les lymphocytes~T qui la reconnaissent sont éteints ou
l'ignorent. Un adjuvant --- l'alun, un lipide, un ligand de \omterm{def:b3:innate-immunity:prr}{récepteur de
type Toll} --- déclenche les récepteurs de motifs, fait mûrir la \omterm{def:b3:innate-immunity:innate}{cellule
dendritique} et fournit le second signal. C'est le point de Janeway~: le
système adaptatif ne répond à un antigène que lorsque le système inné
dit que l'antigène est venu avec un microbe.
\end{solution}

\begin{solution}{exo:b3:innate-immunity:10}
Sans fièvre, $48$ doublements~: $10^{4}\times 2^{48} \approx 3\times
10^{18}$ (une absurdité qui montre que c'est la destruction, et non la
croissance, qui décide de l'issue). Avec la fièvre, $32$ doublements~:
$4\times 10^{13}$ --- un facteur $2^{16} \approx 65\,000$ de bactéries
en moins à affronter pour les \omterm{def:b3:innate-immunity:innate}{neutrophiles}. Coût~:
$3\times 0.12\times 8 = \qty{2.9}{MJ}$ par jour, un gros repas. Cela en
vaut la peine quand le pathogène est ralenti et que l'hôte a des
réserves~; pas quand le pathogène pousse aussi bien à
\qty{40}{\celsius}, ni quand l'hôte est affamé, ni quand la fièvre
elle-même approche des températures dangereuses.
\end{solution}

\begin{solution}{exo:b3:innate-immunity:11}
Fixer le facteur H recrute le régulateur de l'hôte lui-même~: $d$ passe
au-dessus de $a$ et la boucle s'éteint sur le microbe comme sur une
cellule de l'hôte. Cliver le C3b retire la convertase déposée~: encore
une hausse de $d$. Une capsule présente une surface sur laquelle la
convertase se forme mal, ce qui abaisse $a$, et cache aux récepteurs des
phagocytes le C3b déposé. Le moins cher~: une seule protéine de surface
qui fixe le facteur H, que l'hôte fournit gratuitement et qui fait le
travail.
\end{solution}

\begin{solution}{exo:b3:innate-immunity:12}
Au niveau d'une écharde, la dilatation, la fuite et le recrutement
engagent quelques millilitres de vaisseaux et portent la défense là où
il le faut~; les mêmes médiateurs libérés dans toute la circulation
dilatent chaque vaisseau et font fuir chaque capillaire, si bien que la
pression s'effondre et que la coagulation se déclenche partout --- la
physiologie est identique, l'échelle est létale. Le TNF entretient
l'\omterm{def:b3:innate-immunity:inflammation}{inflammation} synoviale de la polyarthrite rhumatoïde, si bien que le
bloquer soulage la maladie~; mais le TNF est aussi ce qui tient les
\omterm{def:b3:innate-immunity:innate}{macrophages} en ordre dans les granulomes qui contiennent les bacilles
tuberculeux dormants, et le bloquer laisse échapper une tuberculose
latente.
\end{solution}

\begin{solution}{pb:b3:innate-immunity:1}
\textbf{1.} Une heure fait $1.5$ doublement~: $10^{4}\times 2^{1.5} \approx
2.8\times 10^{4}$.
\textbf{2.} $2\times 10^{4}$ \omterm{def:b3:innate-immunity:innate}{neutrophiles} tuent jusqu'à $4\times 10^{5}$
bactéries~; les bactéries ne passent que de $2.8\times 10^{4}$ à
$8\times 10^{4}$~: la destruction dépasse la croissance d'un facteur cinq.
\textbf{3.} Heure 2~: $8\times 10^{4}$ avant la destruction, aucune
après. La population culmine vers $8\times 10^{4}$ aux environs de
\qty{2}{h} et est nulle ensuite~; il n'en reste aucune à \qty{6}{h}.
\textbf{4.} Arrivée à \qty{3}{h}~: $10^{4}\times 2^{4.5} = 2.3\times
10^{5}$ alors. Heure 4~: $6.4\times 10^{5} - 4\times 10^{5} = 2.4\times
10^{5}$~; heure 5~: $6.8\times 10^{5} - 4\times 10^{5} = 2.8\times 10^{5}$~;
heure 6~: $7.9\times 10^{5} - 4\times 10^{5} = 3.9\times 10^{5}$ et en
hausse --- les \omterm{def:b3:innate-immunity:innate}{neutrophiles} ne suivent plus, et un abcès se forme. Deux
heures de retard changent une guérison en infection chronique.
\textbf{5.} $8\times 10^{4}/20 = 4000$ \omterm{def:b3:innate-immunity:innate}{neutrophiles} meurent en tuant,
plus les $2\times 10^{4}$ arrivés qui meurent de toute façon dans la
journée. Le pus est fait de \omterm{def:b3:innate-immunity:innate}{neutrophiles} morts, de bactéries mortes et de
tissu liquéfié~; les \omterm{def:b3:innate-immunity:innate}{macrophages} l'évacuent en mangeant les cadavres,
sinon il faut le drainer.
\textbf{6.} Arrivées de $2000$ par heure, destruction de
$4\times 10^{4}$~: heure 2,
$8\times 10^{4} - 4\times 10^{4} = 4\times 10^{4}$~; heure 3, $1.1\times
10^{5} - 4\times 10^{4} = 7\times 10^{4}$~; heure 4, $2\times 10^{5} -
4\times 10^{4} = 1.6\times 10^{5}$ --- croissance sans limite. Le patient
ne peut contenir un inoculum banal, et les \omterm{def:b3:bacteriology:antibiotics}{antibiotiques} doivent donc
tuer dès le premier signe.
\textbf{7.} $B = 12.5(e^{0.08t} - 1)$~: \qty{30}{s}, $125$~; \qty{60}{s},
$1500$~; \qty{120}{s}, $1.8\times 10^{5}$.
\textbf{8.} $e^{0.08t} = 1.6\times 10^{5}$~: $t = 12/0.08 = \qty{150}{s}$.
\textbf{9.} $k_{0}/(d - a) = 10$.
\textbf{10.} $B(120) = 33(e^{3.6} - 1) \approx 1200$, $150$ fois moins.
La capsule achète des minutes --- du temps pour se multiplier, et une
protection jusqu'à l'arrivée des anticorps qui démarreront la voie
classique sur la capsule elle-même.
\textbf{11.} $10^{4}\times 2\times 10^{6} = 2\times 10^{10}$ C3b~: un
millionième du C3 du plasma.
\textbf{12.} Sans C3 les trois voies s'arrêtent à leur étape commune~:
pas d'\omterm{def:b3:innate-immunity:complement}{opsonisation}, pas d'\omterm{def:b3:innate-immunity:complement}{anaphylatoxines}, pas de lyse --- infections
sévères et récidivantes à bactéries encapsulées. Sans C9 seul le pore
terminal est perdu~; l'\omterm{def:b3:innate-immunity:complement}{opsonisation} et l'\omterm{def:b3:innate-immunity:inflammation}{inflammation} fonctionnent, et le
phénotype est fait d'infections récidivantes à \emph{Neisseria}, le seul
genre que les phagocytes contrôlent mal et que la lyse contrôle bien.
\textbf{13.} $2\times 0.24\times 8 = \qty{3.8}{MJ}$, environ \qty{100}{g}
de graisse.
\textbf{14.} Un doublement dans l'heure~: $2\times 10^{4}$ au lieu de
$2.8\times 10^{4}$, un facteur $1.4$.
\textbf{15.} Heure 2~: $4\times 10^{4}$ avant la destruction, aucune
après --- disparues à \qty{2}{h}, comme avant~; l'apport de la fièvre
compte quand l'offre de \omterm{def:b3:innate-immunity:innate}{neutrophiles} est juste, comme aux questions 4 et 6.
\textbf{16.} $199\,\text{mg/L}\times 3\,\text{L} = \qty{0.6}{g}$~; $0.6/
\num{115000} = 5.2\times 10^{-6}$ mol $= 3.1\times 10^{18}$ molécules~;
sur \qty{172800}{s}, $1.8\times 10^{13}$ par seconde.
\textbf{17.} Ils ramènent la consigne vers \qty{37}{\celsius}, le patient
se sent mieux, et les bactéries doublent un peu plus vite~; dans ce
modèle les \omterm{def:b3:innate-immunity:innate}{neutrophiles} l'emportent quand même, et les preuves que les
antipyrétiques aggravent les infections ordinaires sont faibles.
\textbf{18.} Au-dessus de \qty{41}{\celsius} environ les protéines
commencent à se dénaturer, les enzymes défaillent et les neurones
s'emballent (convulsions)~; \qty{42}{\celsius} est létal. La consigne est
tenue en dessous par des cryogènes --- interleukine-10, vasopressine,
glucocorticoïdes --- et par le plafond de l'action des prostaglandines
dans l'hypothalamus.
\textbf{19.} $10^{9}\times 3\times 10^{6} = 3\times 10^{15}$ molécules
$= 5\times 10^{-9}$ mol dans \qty{5}{L}~: \qty{1e-9}{mol/L}.
\textbf{20.} Égale à la concentration saturante~: tous les \omterm{def:b3:innate-immunity:innate}{macrophages} du
corps sont activés au maximum en même temps.
\textbf{21.} $10^{11}\times 1000\times 3600 = 3.6\times 10^{17}$ molécules
$= 6\times 10^{-7}$ mol~; dans \qty{15}{L}, \qty{4e-8}{mol/L} --- quatre
mille fois la concentration active.
\textbf{22.} Chaque artériole se dilate et chaque veinule fuit~: le
volume sanguin stagne dans les vaisseaux élargis et s'échappe dans les
tissus, si bien que la pression tombe~; les reins, mal perfusés, cessent
de filtrer~; et la coagulation déclenchée dans des milliers de
capillaires les bouche et consomme les facteurs de la coagulation.
\textbf{23.} Les \omterm{def:b3:bacteriology:envelope}{Gram négatif} tués libèrent d'un coup tout leur
\omterm{def:b3:bacteriology:envelope}{lipopolysaccharide}, un bolus pour TLR4, si bien que la bouffée de
\omterm{def:b3:innate-immunity:inflammation}{cytokines} peut s'aggraver pendant des heures. L'anti-TNF échoue parce
qu'au moment du choc la cascade a dépassé le TNF --- interleukine-1,
interleukine-6 et coagulation tournent déjà --- et parce que le TNF est
nécessaire pour maîtriser les bactéries~: le problème est le moment, non
la cible.
\textbf{24.} Le \omterm{def:b3:bacteriology:envelope}{lipopolysaccharide} injecté ne fait rien à la souris sans
TLR4, qui survit à des doses létales pour les souris normales~; une
infection vivante à \omterm{def:b3:bacteriology:envelope}{Gram négatif} la tue, parce qu'elle ne peut détecter
tôt les bactéries et monter l'\omterm{def:b3:innate-immunity:inflammation}{inflammation} qui les contient. Le paradoxe
se dissipe~: le détecteur qui cause le choc est le détecteur qui fournit
la défense.
\textbf{25.} Aucune bactérie ne reste à \qty{6}{h} (toutes tuées dès
\qty{2}{h})~; environ $1.8\times 10^{5}$ C3b par bactérie à
\qty{120}{s}~; deux jours de fièvre coûtent \qty{3.8}{MJ}, quelque
\qty{100}{g} de graisse.
\end{solution}
